% Pendahuluan
\section{Pendahuluan}

Ujian Tertulis Berbasis Komputer--Seleksi Nasional Berdasarkan Tes (UTBK-SNBT) 2026 mengukur kemampuan peserta melalui delapan skor subtes. Kedelapan skor itu adalah Penalaran Umum (PU), Pengetahuan dan Pemahaman Umum (PPU), Pemahaman Bacaan dan Menulis (PBM), Pengetahuan Kuantitatif (PK), dua skor Literasi dalam Bahasa Indonesia (LBI), yaitu LBI Saintek dan LBI Soshum, Literasi dalam Bahasa Inggris (LBE), dan Penalaran Matematika (PM).

Statistik deskriptif tiap subtes beredar di media sosial dalam bentuk tabel. Isinya meliputi nilai minimum, kuartil pertama, median, rerata, simpangan baku, kuartil ketiga, dan nilai maksimum. Tabel itu disebut berasal dari materi internal Panitia Seleksi Nasional Penerimaan Mahasiswa Baru (SNPMB), tetapi keasliannya belum diverifikasi secara resmi (lihat Bagian Metode). Oleh karena itu, paper ini menyebut tabel itu \textit{parameter nasional tidak resmi (belum terverifikasi)}. Selain itu, parameter itu tidak mencakup informasi mengenai struktur atau matriks korelasi antarsubtes.

Korelasi menunjukkan seberapa kuat skor suatu subtes berkaitan dengan skor subtes lain. Tanpa informasi korelasi, peserta sulit mengetahui persentil skor kompositnya secara nasional. Skor komposit adalah rerata kedelapan skor subtes, sedangkan persentil menunjukkan persentase peserta yang berskor lebih rendah. Persentil tersebut tidak dapat dihitung hanya dari statistik deskriptif per subtes yang tersedia bagi publik. Hal itu terjadi karena simpangan baku skor komposit, yaitu ukuran sebaran skornya, bergantung langsung pada korelasi antarsubtes pembentuknya.

Data korelasi yang belum tersedia ini juga menyisakan persoalan teoretis dalam psikometri pendidikan. Literatur kemampuan kognitif telah lama mendokumentasikan bahwa skor berbagai tes kognitif cenderung bergerak searah atau berkorelasi positif \autocite{spearman1904,carroll1993,jensen1998}. Pola ini dikenal sebagai \textit{positive manifold} dan lazim dikaitkan dengan satu kemampuan umum yang melandasi berbagai tes, yaitu \textit{g-factor}. Selain \textit{g-factor}, \textit{Carroll's three-stratum theory} \autocite{carroll1993} mengidentifikasi faktor kelompok (\textit{group factors}) pada strata di bawahnya, misalnya kemampuan verbal dan kuantitatif. Dengan demikian, struktur korelasi antarsubtes kognitif berpotensi bersifat hierarkis.

Sejauh penelusuran penulis, belum ada studi empiris atau dokumen resmi yang memublikasikan struktur korelasi antarsubtes UTBK-SNBT, baik pada 2026 maupun pada tahun sebelumnya. Celah literatur ini penting karena UTBK-SNBT menjadi tolok ukur seleksi bagi ratusan ribu peserta setiap tahun. Penelitian ini bertujuan mengisi celah tersebut dengan mengestimasi korelasi antarsubtes UTBK-SNBT 2026. Estimasi memanfaatkan data \textit{crowdsourced}, yaitu skor yang dibagikan secara sukarela oleh banyak peserta di platform media sosial X.

Pembagian skor secara sukarela dapat menyempitkan rentang skor sampel teramati dibandingkan populasi (\textit{restriction of range}), sehingga korelasi sebelum koreksi (\textit{raw sample}) cenderung lebih lemah. Korelasi \textit{raw sample} itu dikoreksi dengan Thorndike Case III, yaitu formula untuk seleksi yang terjadi pada skor komposit, sehingga menghasilkan korelasi populasi terkoreksi. Pendekatan ini memiliki sejumlah keterbatasan metodologis, yang diuraikan secara terbuka pada Bagian Metode dan Keterbatasan Penelitian. Meskipun demikian, hasilnya tetap berupa estimasi empiris awal yang penting selama data populasi resmi tidak tersedia bagi publik.

\subsection{Rumusan Masalah}
Berdasarkan latar belakang di atas, penelitian ini merumuskan empat pertanyaan. Pertama, berapa estimasi rerata korelasi antardelapan subtes UTBK-SNBT 2026 berdasarkan data yang tersedia secara publik di media sosial? Kedua, sejauh mana pengambilan sampel dari media sosial menghasilkan estimasi yang bias, yakni menyimpang secara sistematis dari nilai populasi, dan ke mana arah bias itu? Ketiga, apakah kedelapan subtes membentuk satu \textit{g-factor} tunggal atau struktur berjenjang dengan faktor kelompok? Keempat, bagaimana implikasi estimasi korelasi tersebut terhadap penghitungan persentil skor komposit secara nasional?

\subsection{Tujuan Penelitian}
Sejalan dengan keempat pertanyaan tersebut, penelitian ini bertujuan mengestimasi rerata korelasi antardelapan subtes UTBK-SNBT 2026. Tujuan kedua adalah mengidentifikasi arah dan besaran bias akibat pengambilan sampel dari media sosial. Tujuan ketiga adalah menguji struktur faktor kedelapan subtes melalui \textit{confirmatory factor analysis (CFA)}. Tujuan keempat adalah menganalisis implikasi estimasi korelasi itu bagi penghitungan persentil skor komposit secara nasional.

\subsection{Kontribusi Penelitian}
Penelitian ini memberikan tiga kontribusi utama. Kontribusi substantifnya adalah estimasi empiris awal struktur korelasi antarsubtes UTBK-SNBT 2026, yang hingga kini belum tersedia dalam literatur publik, beserta catatan batasan metodologis yang eksplisit. Kontribusi strukturalnya adalah analisis struktur faktor kedelapan subtes melalui CFA. Kontribusi metodologisnya adalah studi kasus tentang dinamika bias seleksi pada data \textit{crowdsourced} media sosial. Studi kasus ini mengindikasikan bahwa perluasan \textit{framing} (cara menyusun ajakan pengumpulan data) agar lebih inklusif belum tentu mengoreksi bias bila audiens platform sumber data telah tersaring secara sosial. Temuan ini menjadi catatan metodologis penting bagi riset pendidikan berbasis \textit{crowdsourcing} di Indonesia.